%!TEX root = wine_cellar.tex
\documentclass[conference]{IEEEtran}
\IEEEoverridecommandlockouts
% The preceding line is only needed to identify funding in the first footnote. If that is unneeded, please comment it out.
\usepackage{cite}
\usepackage{amsmath,amssymb,amsfonts}
\usepackage{algorithmic}
\usepackage{graphicx}
\usepackage{textcomp}
\usepackage{xcolor}
\usepackage{enumitem}
\usepackage{tabularx}
\def\BibTeX{{\rm B\kern-.05em{\sc i\kern-.025em b}\kern-.08em
    T\kern-.1667em\lower.7ex\hbox{E}\kern-.125emX}}
\begin{document}

\title{AI-Driven Interactive Wine Cellar Management System with Transparent Display Integration\\
}

\author{
\makebox[\textwidth][c]{%
\begin{minipage}[t]{0.245\textwidth}
\centering
\textbf{Naomi K. Fossaberg}\\
\textit{Department of}\
\textit{Computer Science}\\
IT University of Copenhagen\\
Copenhagen, Denmark\\
{\footnotesize nafo@itu.dk}
\end{minipage}%
\hfill
\begin{minipage}[t]{0.245\textwidth}
\centering
\textbf{Enkhbold Myagmarjav}\\
\textit{Department of}\
\textit{Information Systems}\\
Hanyang University\\
Seoul, Republic of Korea\\
{\footnotesize enkheeboldu@hanyang.ac.kr}
\end{minipage}%
\hfill
\begin{minipage}[t]{0.245\textwidth}
\centering
\textbf{Chiang Jun Han}\\
\textit{Department of Electrical}\
\textit{and Electronic Engineering}\\
Nanyang Technological\\
University\\
Singapore\
{\footnotesize jchiang009@e.ntu.edu.sg}
\end{minipage}%
\hfill
\begin{minipage}[t]{0.245\textwidth}
\centering
\textbf{Sungbin Bae}\\
\textit{Department of}\
\textit{Mathematics}\\
Hanyang University\\
Seoul, Republic of Korea\
{\footnotesize charliepro@hanyang.ac.kr}
\end{minipage}%
}
}

\maketitle
\raggedbottom
\begin{abstract}
The premium home appliance sector is increasingly prioritizing artificial intelligence (AI) and user experience (UX) over mere hardware specifications. In the context of luxury wine preservation, existing commercial cellars remain passive climate-control units, while third-party software solutions suffer from high interaction friction and a lack of physical-digital synchronization. To address these bottlenecks, this paper presents an AI-driven, transparent-display-integrated wine cellar management system designed for high-end collections. Our system introduces a robust Vision AI pipeline for automated label recognition and spatial mapping, effectively eliminating manual inventory tracking even for flexibly stacked bottles. Furthermore, we implement a privacy-preserving, on-device Multi-Agent Small Language Model (SLM) architecture. By leveraging a quantized 4-bit base model with dynamic Low-Rank Adaptation (LoRA) swapping on a low-power Neural Processing Unit (NPU), the system generates persona-driven expert dialogues, dynamic food pairings, and personalized daily briefings without violating the stringent thermal and memory constraints of the appliance. Ultimately, this architecture transforms the traditional wine cellar from a passive storage unit into an actively managed, hyper-personalized asset management ecosystem.
\end{abstract}

\begin{IEEEkeywords}
Smart Home Appliance, Edge AI, Vision AI, Small Language Model (SLM), Low-Rank Adaptation (LoRA), Wine Management System, Human-Computer Interaction (HCI)
\end{IEEEkeywords}

\par\smallskip
\noindent\begin{minipage}{\columnwidth}
\centering
\refstepcounter{table}\label{tab:role-assignments}
{\footnotesize TABLE \Roman{table}. \textit{Role Assignments}\par}
\smallskip
\footnotesize
\setlength{\tabcolsep}{3pt}
\renewcommand{\arraystretch}{1.05}
\begin{tabularx}{\columnwidth}{|>{\raggedright\arraybackslash}p{0.24\columnwidth}|>{\raggedright\arraybackslash}p{0.19\columnwidth}|X|}
\hline
\multicolumn{1}{|c|}{Roles} & \multicolumn{1}{c|}{Name} & \multicolumn{1}{c|}{Task description and etc.} \\
\hline{User} & Chiang Jun Han & Describe typical user scenarios and workflows. Identify usability problems and practical needs. Try prototypes and provide feedback.\\
\hline{Customer} & Naomi K. Fossaberg & Define the project's purpose and desired outcomes. Prioritize requirements and features. Evaluate whether the final solution meets the agreed requirements. \\
\hline{Software developer} & Enkhbold Myagmarjav & Design and implement software components. Write and review code. Test functionality and fix bugs. \\
\hline{Development manager} & Sungbin Bae & Plan tasks, milestones, and deadlines. Coordinate collaboration and monitor progress. Identify risks and help resolve blockers.\\
\hline
\end{tabularx}
\end{minipage}\par\smallskip


\section{INTRODUCTION}

\subsection{Motivation}
In recent years, the premium home appliance market has undergone a fundamental paradigm shift, transitioning from pure hardware performance competition toward elevating user experience (UX) through Cloud-based Artificial Intelligence (AI) and the seamless connectivity of the Internet of Things (IoT). High-end luxury product lines, such as LG SIGNATURE, incur substantial hardware manufacturing costs, thereby demanding sophisticated, cloud-backed software capabilities that deliver commensurate value to justify their premium positioning. Concurrently, the global expansion of the fine wine market has driven rapid demand for specialized home wine cellars driven by AI management. However, existing commercial wine cellars remain passive, isolated appliances limited to fundamental condition control (temperature, humidity regulation, anti-vibration, and so on), failing to leverage vast cloud computing resources to satisfy the sophisticated digital management needs of modern fine wine collectors.

\subsection{Problem Statement and Target Audience}
The primary target demographic for luxury wine cellars consists of high-net-worth individuals, predominantly in their 40s to 60s. Although these collectors frequently possess extensive and valuable wine collections, current management paradigms introduce critical usability bottlenecks:
\begin{enumerate}
    \item \textbf{High Interaction Cost and Friction:} Existing mobile-app-based wine management solutions force users to repeatedly handle cold bottles while operating small smartphone screens. This physical separation between the appliance and the management interface induces high cognitive load and interaction friction.
    \item \textbf{Absence of Physical-Digital Synchronization:} Current wine cellars lack real-time synchronization between physical storage and digital databases. Consequently, users are required to perform manual inputs for every bottle movement, leading to severe tracking discrepancies and user abandonment over time.
\end{enumerate}

\subsection{Related Software and Limitations}
An analysis of existing wine management software (e.g., Vivino, CellarTracker) and proprietary cellar tracking apps reveals distinct functional limitations when applied to high-end collections:
\begin{itemize}
    \item \textbf{Absence of Spatial Awareness:} Most commercial applications function merely as flat digital ledgers or list-based inventories. They lack a robust spatial mapping mechanism that correlates a specific bottle's data with its exact physical coordinates (e.g., specific racks and rows) within the cellar, inevitably forcing users to manually search through shelves.
    \item \textbf{Static Data without Environmental Context:} Existing software operates completely independently of the cellar's internal hardware sensors. Consequently, they cannot factor real-time environmental data—such as multi-zone temperature differences—into their aging algorithms or serving readiness calculations.
    \item \textbf{Reactive and Monotone User Experience:} Current tools are inherently reactive, requiring users to explicitly query the database to check a wine's optimal drinking window or food pairing suitability. Furthermore, they rely on static text reviews, lacking the proactive intelligence and dynamic, expert-level engagement necessary to autonomously curate an immersive experience.
\end{itemize}

\subsection{Proposed Solution and Key Contributions}
To overcome these limitations, this paper proposes a cloud-powered, AI-driven wine cellar management system integrated with a transparent display door. Designed specifically for luxury appliances such as the LG SIGNATURE line, the proposed architecture leverages the robust computing infrastructure of centralized datacenters (e.g., LG Datacenter) to deliver uncompromising performance.

The key contributions of this work are summarized as follows:
\begin{itemize}
    \item \textbf{Vision AI and Automated Spatial Mapping:} Integrates a robust vision pipeline to extract label and vintage data. Upon insertion, an automated placement mechanism registers the bottle's location. This extends beyond rigid 2D grids, possessing the capability to accurately map bottles even when they are flexibly stacked on top of one another without discrete rack partitions.
    \item \textbf{Cloud-Backed Multi-Agent LLM Architecture:} Replaces resource-constrained edge computing with a high-performance cloud backend (LG Datacenter). This paradigm shift allows the system to deploy massive, state-of-the-art Large Language Models (LLMs) without imposing any thermal, processing, or memory constraints on the cellar appliance itself. The cloud architecture seamlessly orchestrates multi-agent dialogues among critic personas, ensuring hyper-personalized, expert-level interactive experiences.
    \item \textbf{AI-Driven UX Enhancement:} Transcending basic touch interactions, the cloud-integrated AI engine drastically elevates the overall user experience. By analyzing multidimensional real-time data, the system proactively offers a broad spectrum of personalized management features—such as optimal drinking window alerts and dynamic food pairings—transforming passive wine storage into an actively managed ecosystem.
    \item \textbf{Multi-Device Distributed Scalability:} Establishes a synchronized system architecture that seamlessly unifies multiple physical cellar units under a single user profile. Leveraging cloud centralization, it guarantees real-time synchronization between the mobile application and the appliance, eradicating tracking discrepancies.
\end{itemize}

\section{REQUIREMENTS SPECIFICATION}

\subsection{Functional Requirements}

\subsubsection{User Management}
\begin{itemize}
    \item The system shall allow a user to create an account with the following mandatory information: Username, Email, and Password.
    \item The system shall allow a registered user to log in and out.
    \item The system shall allow an authenticated user to edit their profile information.
\end{itemize}

\subsubsection{Inventory Mapping}
\begin{itemize}
    \item The system shall automatically detect wine bottles placed inside the wine cellar using computer vision.
    \item The system shall extract wine label information from detected bottles using AI.
    \item The system shall automatically register newly detected bottles in the inventory.
    \item The system shall automatically associate each detected bottle with a physical location in the wine cellar.
    \item The system shall automatically update a bottle's location when it is moved within the wine cellar.
    \item The system shall assign a unique identifier to each physical wine bottle to prevent duplicate inventory records.
    \item The system shall allow an authenticated user to review and correct AI-extracted wine information before confirming registration.
    \item The system shall automatically detect stacked wine bottles.
    \item The system shall automatically remove a bottle from the inventory when it is removed from the wine cellar.
    \item The system shall allow an authenticated user to manually add a wine bottle to the inventory.
    \item The system shall allow an authenticated user to manually modify bottle information.
    \item The system shall allow an authenticated user to manually remove a wine bottle from the inventory.
\end{itemize}

\subsubsection{Inventory Search \& Visualisation}
\begin{itemize}
    \item The system shall allow a user to search for a bottle by wine name, winery, region, vintage or grape variety.
    \item The system shall display the physical location of a selected bottle within the wine cellar.
    \item The system shall display a visual map of bottles stored in the wine cellar.
\end{itemize}

\subsubsection{Environmental Monitoring}
\begin{itemize}
    \item The system shall collect environmental data from the wine cellar sensors.
    \item The system shall store historic environmental measurements from the wine cellar.
    \item The system shall associate environmental conditions with individual bottle locations.
    \item The system shall notify users when temperature or humidity exceeds user-defined thresholds for a specified duration.
    \item The system shall notify users when sensors become disconnected or environmental data is outdated.
\end{itemize}

\subsubsection{AI Recommendations}
\begin{itemize}
    \item The system shall allow a user to input data regarding a meal and receive wine pairing recommendations from the user's current inventory.
    \item The system shall allow a user to select a specific wine from the inventory and generate complementary food pairing suggestions, including culinary recipes or food delivery menu categories.
    \item The system shall recommend bottles that have reached their optimal drinking window.
    \item The system shall notify a user when wines are approaching their optimal drinking window.
    \item The system shall generate personalised recommendations based on historical consumption patterns.
    \item The system shall allow a user to input specific constraints, such as a target budget, flavor profile, or event theme, to request a curated wine list.
    \item The system shall automatically generate and display a customized wine list that satisfies the user-defined parameters, utilizing both the existing cellar inventory and external databases for potential new purchases.
    \item The system shall recommend suitable storage locations for newly added bottles.
    \item The system shall recommend wines via natural language queries processed by the cloud-based LLM.
\end{itemize}

\subsubsection{Personalized Daily Wine Briefing}
\begin{itemize}
    \item The system shall aggregate daily global wine news, vintage reports, and market trends from external web sources and APIs.
    \item The system shall personalize the aggregated news content by cross-referencing it with the user's specific cellar inventory, highlighting articles relevant to their owned regions, grapes, or producers.
    \item The system shall automatically curate and display a "Daily Wine Briefing" on the cellar's front display, including inventory-specific insights (e.g., wines reaching their peak, humidity stability).
    \item The system shall send a summarized version of the daily briefing as a push notification to the user's synchronized mobile application at a user-defined time.
\end{itemize}

\subsubsection{Real-Time Market Value Integration}
\begin{itemize}
    \item The system shall integrate with external global wine market APIs (e.g., Wine-Searcher) to retrieve real-time pricing data for the specific vintages stored in the cellar.
    \item The system shall calculate and visualize the real-time aggregated market value of the user's entire wine collection on the front display.
    \item The system shall provide graphical visualizations of historical price trends and value appreciation for individual bottles to assist the user in asset management.
\end{itemize}

\subsubsection{Multi-Cellar Management}
\begin{itemize}
    \item The system shall allow a user to manage multiple wine cellar units and wine cellars under a single account.
    \item The system shall synchronise inventory information across connected wine cellars.
    \item The system shall allow a user to search across all connected wine cellars.
\end{itemize}

\subsubsection{Cloud-Powered Multi-Agent LLM Dialog System}
\begin{itemize}
    \item The system shall integrate a cloud-hosted Multi-Agent Large Language Model (LLM) framework, configured with multiple distinct personas reflecting the evaluation tendencies and grading philosophies of renowned wine critics (e.g., Robert Parker, Luca Maroni).
    \item The system shall automatically send an API request to the cloud backend to trigger a simulated multi-agent dialogue whenever a new wine is registered to the cellar inventory or actively selected by the user via the touch display.
    \item The system shall retrieve and visualize interactive, chat-style debates on the front display in real-time, showcasing diverse analytical perspectives, tasting notes, and aging predictions for the selected wine.
    \item The system shall allow the cloud-based AI to continuously fine-tune its conversational outputs and recommendations by securely analyzing the user's synchronized cellar database, including inventory history, consumption patterns, and individual preferences.
    \item \textbf{Cloud Base Model Operation:} The system shall utilize a centralized, high-parameter Large Language Model (LLM) hosted on the data center as the foundational reasoning engine, removing the need for aggressive quantization.
    \item \textbf{Server-Side Dynamic LoRA Swapping:} The system shall dynamically load and unload specific Low-Rank Adaptation (LoRA) weight matrices corresponding to the active persona on the cloud server during multi-agent dialogue generation, maximizing GPU VRAM efficiency while ensuring perfect persona emulation.
    \item \textbf{Continuous Persona Expansion:} The system shall allow the seamless addition of new critic or user-specific personas entirely via server-side LoRA updates, requiring no local software patches or model downloads on the cellar appliance.
\end{itemize}

\subsubsection{Mobile Application Integration}
\begin{itemize}
    \item The system shall provide a companion mobile application that synchronizes real-time inventory and cellar status with the user's account.
    \item The system shall allow a user to remotely browse their wine inventory, search for bottles, and view AI recommendations via the mobile application.
    \item The system shall push real-time notifications to the mobile application regarding critical environmental alerts (e.g., temperature fluctuations) and optimal drinking windows.
    \item The system shall allow a user to remotely adjust the temperature and humidity settings of individual cellar zones through the mobile application.
    \item The system shall allow a user to manually add, edit, or remove wine inventory data via the mobile application when away from the physical cellar.
\end{itemize}

\subsection{Non-Functional Requirements}

\subsubsection{Performance}
\begin{itemize}
    \item The system shall return inventory search results within 2 seconds for 95\% of requests.
    \item The system shall update the displayed inventory within 10 seconds after a bottle had been added or removed.
    \item The system shall process API requests to the cloud-based LLM and generate the initial multi-agent dialogue response within 1.5 seconds, utilizing token streaming to ensure a seamless user experience despite network latency.
\end{itemize}

\subsubsection{Resource \& Hardware Constraints}
\begin{itemize}
    \item The cellar's internal computing module shall function strictly as a thin client, capping local memory usage (RAM) at 1 GB and offloading all intensive Vision AI and LLM inferences to the cloud.
    \item The local computing processes shall be strictly low-power to ensure zero thermal interference with the cellar's precise climate control mechanisms.
\end{itemize}

\subsubsection{Reliability \& Availability}
\begin{itemize}
    \item The system shall preserve the submitted data locally so that it remains after a temporary application restart.
    \item The system shall implement an offline fallback mechanism, retaining basic inventory viewing and temperature control capabilities when the Wi-Fi connection to the cloud is temporarily lost.
    \item The system shall automatically back up all inventory and configuration data to the cloud servers and support seamless recovery following an unexpected shutdown or hardware replacement.
\end{itemize}

\subsubsection{Security \& Privacy}
\begin{itemize}
    \item The system shall require secure authentication (e.g., OAuth 2.0) before a user can access inventory management functions across devices.
    \item The system shall securely hash and store user passwords on the cloud backend.
    \item The system shall enforce End-to-End (E2E) encryption (TLS 1.3) for all data transmissions between the cellar appliance, the mobile application, and the LG Datacenter, ensuring personal consumption data and AI interactions cannot be intercepted.
    \item The system shall strictly isolate user databases on the cloud server to prevent cross-account data leaks.
    \item The system shall be fully compliant with data protection regulations such as GDPR, offering users the right to export or permanently delete their data from the cloud.
\end{itemize}

\subsubsection{Maintainability}
\begin{itemize}
    \item The system architecture shall support the server-side replacement, scaling, or upgrading of the backend LLMs without requiring firmware updates or modifications to the user interface on the cellar appliance.
    \item The system architecture shall cleanly separate the frontend UI components from the backend API services to allow independent development and deployment cycles.
\end{itemize}

\subsubsection{Usability}
\begin{itemize}
    \item The system shall provide an intuitive interface allowing new users to navigate core features without consulting external documentation.
    \item The system shall clearly indicate the current network connection status and authentication state on the main display.
    \item The system shall preserve unsaved user input locally during sudden network drops, preventing data loss before synchronizing with the cloud.
\end{itemize}

\section{Development Environment}

\subsection{Choice of Software Development Platform}

\subsubsection{Development Operating Systems}
\begin{itemize}
    \item \textbf{Windows 11 with WSL2 (Ubuntu 22.04.5 LTS):} Windows 11 is selected as the primary development operating system for the backend and the integrated display because it runs all of the main project tools, including Visual Studio Code, Android Studio, Docker, and Git. Windows Subsystem for Linux 2 (WSL2) provides a full Ubuntu environment alongside Windows without a separate virtual machine, and is used for Python development, dependency management, and Docker-based workflows. Because the AWS EC2 server also runs Ubuntu, developing in WSL2 reduces differences between local and cloud environments.
    \item \textbf{macOS Golden Gate (27.2):} macOS is selected for the native iOS application because Apple's development tools, including Xcode, the Swift compiler, SwiftUI, and the iOS Simulator, are only available on macOS. Building and testing an iOS application therefore requires at least one macOS machine within the team. Its Unix-based architecture also provides a terminal environment similar to Linux, allowing the same command-line tools and Git workflow used for the other components to be applied without modification.
\end{itemize}

\subsubsection{Target Operating Systems}
\begin{itemize}
    \item \textbf{Android (Android 13):} Android is selected as one of the target platforms for the companion mobile application, which gives users remote access to the system when they are away from the appliance. As the platform with the largest global smartphone market share, it is essential for reaching most potential users. The application allows users to sign in, browse and manage their inventory, monitor temperature and humidity, receive notifications, and view AI recommendations. It is developed natively with Kotlin and Jetpack Compose and stores all data through the shared FastAPI backend.
    \item \textbf{iOS (iOS 17):} iOS is selected as the second target mobile platform to provide the same core functionality for iPhone users. The iOS application is a separate native implementation written in Swift and SwiftUI, allowing it to follow Apple's interface guidelines rather than imitating the Android design. Although the two applications are implemented independently, both use the same FastAPI backend and API contracts. As a result, business logic and data remain centralized, and an inventory change made on one device is reflected consistently on all other clients.
    \item \textbf{webOS Open Source Edition (2.28.0):} webOS Open Source Edition (OSE) is selected as the intended target operating system for the integrated transparent display because it is designed for connected and embedded devices and runs applications built with standard web technologies. The display can therefore be implemented as a React and TypeScript web application, developed and demonstrated in a standard browser, and hosted on Vercel during development without physical wine cellar hardware. Because the display communicates with the backend only through its APIs, it can later be moved to webOS without changes to the backend services.
\end{itemize}

\subsubsection{Programming Languages}
\begin{itemize}
    \item \textbf{Python (3.12):} Python is selected as the primary backend language because of its readable syntax and extensive libraries for web services, data processing, computer vision, and artificial intelligence. It implements the FastAPI backend, inventory and cellar management logic, environmental sensor data processing, bottle localization, and requests to external LLM services. Using one language for both application services and data processing reduces the complexity of integrating image recognition and AI-powered recommendations into the backend. The backend runs in a Python 3.12 Docker container on AWS EC2, so all developers and the server use the same interpreter version.
    \item \textbf{TypeScript (5.9.3):} TypeScript is selected for the integrated display because it extends JavaScript with static typing that detects many errors during development rather than at runtime. The display processes structured data from the backend, including wine metadata, bottle identifiers, storage positions, sensor readings, and AI responses. Defining these structures as TypeScript interfaces that mirror the API responses allows integration errors to be caught at compile time. TypeScript is used with React to implement the visual cellar map, inventory search, bottle details, and AI conversation interface.
    \item \textbf{Kotlin (2.4.21):} Kotlin is selected for the native Android application because it is Google's recommended language for Android development and integrates directly with Android Studio, the Android SDK, and Jetpack Compose. Its concise syntax reduces boilerplate code compared to Java, and its null-safety system prevents runtime errors caused by missing values in backend responses. Kotlin coroutines allow network requests to run without blocking the user interface. It is used to implement authentication, navigation, inventory management, environmental monitoring, and communication with the backend.
    \item \textbf{Swift (6.4):} Swift is selected for the native iOS application because it is Apple's primary programming language and integrates directly with SwiftUI and the iOS SDK. Its strong type system and optional types make the handling of missing or invalid data explicit, reducing errors when decoding backend responses. Its async/await concurrency model allows network requests to be written clearly without blocking the user interface. Swift is used to implement navigation, authentication, inventory operations, environmental monitoring, and AI-related interactions with the FastAPI backend.
    \item \textbf{SQL (PostgreSQL Dialect):} SQL is selected for managing the relational data at the core of the system. The database maintains relationships between users, wine products, individual physical bottles, cellar units, storage positions, and environmental measurements; many bottles can refer to the same wine product while each keeps its own identifier and location. Primary keys, foreign keys, and constraints enforce these relationships and prevent inconsistent records. Most queries are issued through SQLAlchemy, while SQL is used directly for schema definitions, migrations, and inspecting data during development.
\end{itemize}

\subsubsection{Cloud Infrastructure}
\begin{itemize}
    \item \textbf{Amazon Web Services (AWS) --- Amazon EC2:} Amazon EC2 is selected as the main cloud hosting environment because it provides configurable virtual machines with full control over the operating system, networking, and deployed software. A single Linux t3.small instance runs the FastAPI backend and the PostgreSQL database in separate Docker containers coordinated by Docker Compose. Because language model inference is handled by external services, no GPU is required, and the 2~vCPU, 2~GB configuration is sufficient for development and demonstration at low cost. This provides one central backend for the integrated display and both mobile applications.
    \item \textbf{Vercel:} Vercel is selected to host the React-based integrated display because it specializes in deploying frontend web applications and supports Vite projects without additional configuration. Its GitHub integration automatically builds and deploys the application when changes are pushed, and creates preview deployments for branches so that interface changes can be reviewed before merging. The deployed display retrieves inventory, cellar status, environmental readings, and AI results from the FastAPI backend. Hosting the frontend separately allows it to be developed and deployed independently of the backend.
    \item \textbf{Hugging Face Inference Providers (Testing):} Hugging Face's free inference tier is selected for development and testing because it provides API access to open-weight language models without cost or a dedicated GPU server. The team uses an instruction-tuned open-weight model, Qwen2.5-7B-Instruct, to develop and test prompts, critic personas, and the multi-agent dialogue flow for wine recommendations and food pairings. Because the free tier includes only limited monthly usage credits and offers no guaranteed response times, it is used during development only and not in the deployed system.
    \item \textbf{OpenAI API (Deployment):} The OpenAI API is selected for the deployed system because it offers reliable, low-latency access to high-quality language models with usage-based pricing. The team plans to use GPT-5 mini, a cost-efficient model suited to recommendation, food pairing, and critic dialogue tasks. For each request, the FastAPI backend retrieves relevant inventory data and user preferences from PostgreSQL, builds persona-specific prompts, and combines the critic responses into a stored conversation. Token streaming shortens the time until the first response appears. Both providers are accessed through a common backend interface, so switching from testing to deployment requires only a configuration change.
\end{itemize}

\subsubsection{Development Tools}
\begin{itemize}
    \item \textbf{Visual Studio Code (1.144.1):} Visual Studio Code is selected as the primary code editor because it is lightweight, free, and supports Python, TypeScript, Docker, and Git through an extensive ecosystem of extensions. It provides code completion, linting, and an integrated debugger for both the FastAPI backend and the React display. Its WSL integration allows files inside the Ubuntu environment to be edited directly, and its integrated terminal allows developers to run backend services and execute Docker and Git commands without leaving the editor.
    \item \textbf{Xcode (27.2):} Xcode is selected as the development environment for the native iOS application because it is Apple's official IDE and includes the iOS SDK, the Swift compiler, and debugging and profiling tools. SwiftUI previews allow interface components to be inspected while they are being written. The iOS Simulator is used to test the application across different iPhone models and screen sizes, while testing on physical devices validates behavior and network communication outside the simulated environment.
    \item \textbf{Android Studio Rabbit 2 (2026.2.2):} Android Studio is selected for the native Android application because it is Google's official IDE and provides integrated support for Kotlin, Jetpack Compose, Gradle, and the Android SDK. Compose previews allow screens to be inspected without launching the full application, and built-in profiling tools help identify performance issues. The Android Emulator is used to test the application across different device configurations and Android versions, while physical-device testing verifies behavior under practical network and hardware conditions.
    \item \textbf{Docker Engine (29.9.2) and Docker Compose (5.7.0):} Docker is selected to package the backend services and their dependencies into containers, ensuring that the same configuration runs on every developer's machine and on the AWS EC2 server. The FastAPI backend and PostgreSQL database run in separate containers, so each can be updated or restarted independently. Docker Compose defines both containers, their internal network, environment variables, and persistent storage volumes in a single file. Any team member can start the complete backend locally with one command, and the same configuration is used for deployment.
    \item \textbf{Git (2.34.1) and GitHub:} Git is selected as the version control system because several components are developed in parallel while sharing one source code history. Each feature is developed on a separate branch and integrated through merging, while the commit history allows previous versions to be restored when necessary. GitHub hosts the central repository and adds pull requests, code reviews, and issue tracking. The repository contains separate directories for the backend, integrated display, Android application, iOS application, and documentation, and its Vercel integration automatically deploys the integrated display.
    \item \textbf{Figma:} Figma is selected for interface design and prototyping because it is a browser-based tool that allows all team members to edit the same designs in real time. The project requires a consistent user experience across the integrated display and two native mobile applications, and Figma's shared components and styles allow common elements such as colors, typography, and wine information cards to be defined once and reused. Interactive prototypes allow navigation and user flows to be evaluated before implementation, reducing interface changes during development.
    \item \textbf{Draw.io (diagrams.net):} Draw.io is selected to create technical diagrams because it is free, runs in the browser, and stores diagrams in a format that can be tracked in the Git repository. It is used to represent the system's components, including the wine cellar controller, integrated display, FastAPI backend, PostgreSQL database, AI services, and mobile applications, along with the communication paths between them. It is also used to prepare data flow diagrams and process flowcharts for the specification and design documents.
    \item \textbf{LaTeX:} LaTeX is selected for the project documentation because it produces consistently formatted technical documents and handles tables, figures, equations, and references reliably. The team uses the IEEE conference template, as required for the course, to organize the introduction, requirements, development environment, specifications, and subsequent documentation. Because LaTeX documents are plain text, the source files are stored in the GitHub repository alongside the software, allowing documentation changes to be tracked and reviewed in the same way as source code.
\end{itemize}

\subsubsection{AI-Assisted Development Tools}
\begin{itemize}
    \item \textbf{Codex CLI (0.162.1; ChatGPT Plus; Astra 6.1):} Codex CLI is a command-line coding agent developed by OpenAI, used in this project with the Astra 6.1 model through a ChatGPT Plus subscription. It works directly within the project's source files and terminal, allowing it to read existing code, propose changes, and run tests and builds. The team uses Codex mainly for implementing individual features, including backend API endpoints, React components, and database operations, as well as for debugging and writing automated tests.
    \item \textbf{Claude Code (2.1.309; Claude Pro; Sonnet 5.5):} Claude Code is a command-line coding agent developed by Anthropic, used in this project with the Claude Sonnet 5.5 model through a Claude Pro subscription. Like Codex, it works directly with the repository and development tools, but the team uses it mainly for tasks that span several components. These include refactoring code across modules, reviewing implementation approaches against the planned architecture, and checking consistency between the backend API and the data structures used by the client applications.
\end{itemize}

\subsubsection{Computer Resources}
Table~\ref{tab:computer-resources} lists the computer resources used for development and deployment.

\par\smallskip
\noindent\begin{minipage}{\columnwidth}
\centering
\refstepcounter{table}\label{tab:computer-resources}
{\footnotesize TABLE \Roman{table}. \textit{Computer Resources}\par}
\smallskip
\footnotesize
\setlength{\tabcolsep}{3pt}
\renewcommand{\arraystretch}{1.05}
\begin{tabularx}{\columnwidth}{|>{\raggedright\arraybackslash}p{0.24\columnwidth}|X|>{\raggedright\arraybackslash}p{0.16\columnwidth}|>{\raggedright\arraybackslash}p{0.12\columnwidth}|}
\hline
\multicolumn{1}{|c|}{Resource} & \multicolumn{1}{c|}{Device / OS} & \multicolumn{1}{c|}{CPU} & \multicolumn{1}{c|}{RAM} \\
\hline Naomi K. Fossaberg & & & \\
\hline Enkhbold Myagmarjav & ASUS Zenbook 14 (UX3405MA), Windows 11 Home (Build 26300) & Intel Core Ultra 7 155H (16 cores, 22 threads) & 16 GB \\
\hline Chiang Jun Han & & & \\
\hline Sungbin Bae & & & \\
\hline AWS EC2 (t3.small) & Ubuntu 22.04 LTS, 20 GB gp3 storage, Seoul region & 2 vCPU & 2 GB \\
\hline
\end{tabularx}
\end{minipage}\par\smallskip

\subsubsection{Cost Estimation}
\mbox{}\par\smallskip
\noindent\begin{minipage}{\columnwidth}
\centering
\refstepcounter{table}\label{tab:cost-estimation}
{\footnotesize TABLE \Roman{table}. \textit{Estimated Development Cost}\par}
\smallskip
\footnotesize
\setlength{\tabcolsep}{3pt}
\renewcommand{\arraystretch}{1.05}
\begin{tabularx}{\columnwidth}{|>{\raggedright\arraybackslash}X|>{\raggedright\arraybackslash}p{0.24\columnwidth}|}
\hline
\multicolumn{1}{|c|}{Item} & \multicolumn{1}{c|}{Cost (USD)} \\
\hline Development laptops & \$3,500 \\
\hline AWS EC2 t3.small, 20 GB storage, and public IPv4 address & \$25--30 \\
\hline Vercel (Hobby plan) & \$0 \\
\hline Hugging Face Inference (free tier, testing) & \$0 \\
\hline OpenAI API usage (GPT-5 mini, deployment) & \$10--20 \\
\hline ChatGPT Plus (Codex CLI) & \$20 \\
\hline Claude Pro (Claude Code) & \$20 \\
\hline VS Code, Xcode, Android Studio, Docker, Git, GitHub, Figma, Draw.io, and LaTeX (free tools/plans) & \$0 \\
\hline \textbf{Total} & \textbf{\$3,575--3,590} \\
\hline
\end{tabularx}
\end{minipage}\par\smallskip

\noindent\textbf{Cost Breakdown Notes:}
\begin{itemize}
    \item \textbf{Hardware:} Development laptops for the team, including a macOS machine required for iOS development.
    \item \textbf{Cloud Infrastructure:} An AWS EC2 t3.small instance in the Seoul region hosts the backend and database, including storage and a public IPv4 address. Vercel's free Hobby plan hosts the integrated display.
    \item \textbf{AI Services:} Hugging Face's free inference tier is used for testing, while the OpenAI API (GPT-5 mini) is billed by usage in deployment.
    \item \textbf{Development Tools:} Most tools used in the project, including VS Code, Xcode, Android Studio, Docker, Git, GitHub, Figma, Draw.io, and LaTeX, are free or open source.
    \item \textbf{AI Coding Assistants:} ChatGPT Plus (Codex CLI) and Claude Pro (Claude Code) subscriptions are used throughout development.\end{itemize}

\subsection{Software in Use}

\subsubsection{Frontend Frameworks and Build Tools}
\begin{itemize}
    \item \textbf{React (19.3.1):} React is selected as the frontend library for the integrated transparent display because its component-based architecture allows the interface to be built from small, reusable elements that manage their own state. The visual cellar map, bottle information, environmental measurements, and AI conversations are each implemented as separate components. When new data is received from the backend, React updates only the affected parts of the interface without reloading the page, keeping the display responsive. Its large ecosystem and first-class TypeScript support also simplify development and maintenance.
    \item \textbf{Vite (7.3.7):} Vite is selected as the development server and build tool for the integrated display because it provides a fast development workflow and produces optimized assets for deployment. During development, its hot module replacement shows changes to a component in the running interface almost immediately without losing the current application state. For deployment, it bundles and minifies the application into static files served by Vercel. Using Vite keeps the frontend build process simple and independent of the backend and database infrastructure.
\end{itemize}

\subsubsection{Backend Frameworks}
\begin{itemize}
    \item \textbf{FastAPI (0.162.0):} FastAPI is selected as the primary backend framework because it is a modern, high-performance Python framework designed for building APIs. It uses Python type hints to validate requests and serialize responses automatically, catching malformed data before it reaches the application logic. It supports asynchronous request handling and WebSocket connections natively, and automatically generates OpenAPI documentation that serves as a shared API reference for all client applications. FastAPI implements authentication, inventory management, bottle location updates, cellar configuration, environmental monitoring, and AI service communication, so all clients share the same business logic.
\end{itemize}

\subsubsection{Mobile Application Frameworks}
\begin{itemize}
    \item \textbf{SwiftUI (iOS 27.2 SDK):} SwiftUI is selected for the native iOS application because it is Apple's modern declarative framework for building user interfaces with Swift. Developers describe how the interface should appear for a given state, and SwiftUI updates the screen automatically when that state changes. This allows inventory, cellar, and environmental data received from the backend to be reflected immediately in the corresponding views. SwiftUI is used to build the authentication, inventory, bottle detail, environmental monitoring, and AI recommendation screens, with built-in navigation and accessibility support.
    \item \textbf{Jetpack Compose (UI/runtime 1.12.1):} Jetpack Compose is selected for the native Android application because it is Android's modern declarative UI toolkit and is fully integrated with Kotlin and Android Studio. Like SwiftUI, Compose describes the interface for a given state and automatically updates the affected components when that state changes, for example when new inventory information is retrieved from the backend. This removes the need for XML layouts and much manual view-update code. Because both mobile applications use a declarative approach, the team can follow a similar screen structure on both platforms.
\end{itemize}

\subsubsection{Database}
\begin{itemize}
    \item \textbf{PostgreSQL (17):} PostgreSQL is selected as the primary database because it is a mature, open-source relational database known for its reliability and data integrity. Its transactions and relational constraints ensure that operations such as registering, moving, or removing a bottle are either fully completed or not applied at all. Indexing supports fast inventory searches by wine name, winery, region, vintage, or grape variety, while JSON data types allow semi-structured data such as AI conversation history to be stored. PostgreSQL runs in a Docker container on AWS EC2 with persistent storage, providing one central data source for all clients.
    \item \textbf{SQLAlchemy (2.0):} SQLAlchemy is selected as the Python Object-Relational Mapper (ORM) because it is the most widely used database toolkit in the Python ecosystem and integrates well with FastAPI. It allows tables and relationships, such as the one-to-many relationship between a wine product and its physical bottles, to be defined as Python classes and translates operations on these objects into SQL. It also manages database sessions and transactions. Using an ORM reduces repetitive SQL code and keeps database operations separate from FastAPI route handlers.
\end{itemize}

\subsubsection{Communication and Integration}
\begin{itemize}
    \item \textbf{REST API over HTTP/HTTPS with JSON:} RESTful APIs are selected as the primary communication mechanism because they provide a simple, stateless interface supported on every platform in the project. Resources such as users, wines, bottles, cellars, and sensor readings are exposed through endpoints that use standard HTTP methods (GET, POST, PUT, PATCH, and DELETE). HTTP is used during local development, while HTTPS secures the deployed backend through Transport Layer Security (TLS). Data is exchanged in JSON, which is supported in Python, TypeScript, Kotlin, and Swift, allowing all clients to use a single backend.
    \item \textbf{WebSocket:} WebSocket is selected to support persistent, bidirectional communication for features that require continuous updates. Unlike REST requests, which require the client to ask for new data, a WebSocket connection remains open, allowing the backend to send updates to clients as soon as changes occur. It is used to transmit changes in inventory, bottle locations, environmental measurements, and cellar connection status, so that an action on one device is reflected on all other clients without repeated polling. REST APIs remain responsible for standard request-response operations.
\end{itemize}

\subsection{Task Distribution}

\par\smallskip
\noindent\begin{minipage}{\columnwidth}
\centering
\refstepcounter{table}\label{tab:task-distribution}
{\footnotesize TABLE \Roman{table}. \textit{Task Distribution}\par}
\smallskip
\footnotesize
\setlength{\tabcolsep}{3pt}
\renewcommand{\arraystretch}{1.05}
\begin{tabularx}{\columnwidth}{|>{\raggedright\arraybackslash}p{0.24\columnwidth}|>{\raggedright\arraybackslash}p{0.19\columnwidth}|X|}
\hline
\multicolumn{1}{|c|}{Role} & \multicolumn{1}{c|}{Name} & \multicolumn{1}{c|}{Task Description} \\
\hline Backend \& Cloud Developer & & \\[1em]
\hline Integrated Display Developer & & \\[1em]
\hline Mobile Application Developer & & \\[1em]
\hline AI Integration Developer & & \\[1em]
\hline
\end{tabularx}
\end{minipage}\par\smallskip


\section{Specifications}

\end{document}
